\chapter{Discussão}\label{cap:discussao}

Os resultados não sustentam a hipótese central. Com a linha de base de uma rodada, o detector principal marcou parte da operação normal medida depois de mudanças na montagem e não marcou assinaturas de desbalanceamento do tamanho que a norma considera significativo. A menor mudança detectável em $d_1$ passou de \SI{0,7}{\milli\meter\per\second} em três dos cinco ensaios avaliados, e, na base pública, depois da remontagem do suporte, o detector marcou \hQuatroFracaoZero~\% das janelas sem desbalanceamento. As análises posteriores mostraram dois limites do próprio detector. O falso alarme estimado com poucas janelas de treino variou com a semente do gerador de números aleatórios, a ponto de H1 ser corroborada com parte das sementes, e o escore do Isolation Forest tem um teto que impediu a marcação da assinatura em três das quatro máquinas. A cadeia de aquisição entregou completas todas as janelas dos trechos analisados, e as perdas do enlace ficaram fora deles (seção~\ref{sec:disc-sistema}).

\section{Falso alarme e estabilidade da linha de base}\label{sec:disc-falso-alarme}

O detector principal marcou operação normal medida depois de mudanças na montagem. Na calandra, treinado na primeira rodada, ele marcou \hUmCalandraIFDireto~\% das janelas da segunda, com intervalo de \hUmCalandraIFDiretoInf~\% a \hUmCalandraIFDiretoSup~\% (seção~\ref{sec:res-h1}). Entre as rodadas, mudaram ao mesmo tempo a fixação, a alimentação, o intervalo entre ciclos e a temperatura do sensor, de \campoTempCalandraUm\ para \campoTempCalandraDois~\si{\celsius}, e os dados não indicam qual dessas mudanças produziu a marcação. Na base pública, o falso alarme em 0E foi de \hQuatroFracaoZero~\%, depois da remontagem do suporte do desbalanceamento, e, de 0D para 0E, a mediana de $d_1$ passou de \hQuatroDumMedianaTreino\ para \hQuatroDumMedianaZero\ na unidade dos canais (seção~\ref{sec:res-h4}). Nos dois casos, o treino e o teste foram medidos na mesma máquina sem falha, em montagens diferentes, e o detector não distingue essa diferença de uma falha.

Com poucas janelas de treino, o falso alarme estimado dependeu da semente do detector e do próprio treino. Na calandra, com as sementes de 0 a \num{99}, a fração marcada teve mediana de \apTresSementeCalandraMed~\%, abaixo dos \hUmCalandraIFDireto~\% da semente registrada, e H1 seria corroborada com \apTresVereditoUmCorroborada\ das \apTresSementes\ sementes (seção~\ref{sec:res-ap3}). Com o treino reamostrado e o limiar refeito, o intervalo da calandra foi de \apTresTreinoInfCalandra~\% a \apTresTreinoSupCalandra~\%. Na base pública, com \hQuatroNTreino\ janelas de treino, \SI{95}{\percent} das sementes deram falso alarme em 0E de \apTresSementeNivelZeroPInf~\% a \apTresSementeNivelZeroPSup~\%. Na análise de valores discrepantes, o limiar depende do número de observações e do número de dimensões \cite[p.~650]{worden2000}, e o limiar deste trabalho foi estimado com \campoJanelasGuilhotinaUm\ a \campoJanelasCalandraUm\ janelas e seis descritores. O veredito indeterminado de H1 reflete, portanto, a variação do próprio detector treinado com \campoJanelasCalandraUm\ janelas, além do falso alarme entre rodadas, e uma estimativa que varia de \apTresTreinoInfCalandra~\% a \apTresTreinoSupCalandra~\% não permite fixar um limiar de operação.

Na guilhotina, o falso alarme de \hUmGuilhotinaIFDireto~\% se repetiu com as \apTresSementes\ sementes, e o mesmo modelo marcou \condGuilhotinaParado~\% das janelas com o motor desligado (seção~\ref{sec:res-condicao}). Com \campoJanelasGuilhotinaUm\ janelas de treino, o detector da guilhotina quase não marca janela alguma, e o falso alarme nulo não mede a estabilidade da linha de base dessa máquina.

O escore do Isolation Forest tem um teto fora da faixa do treino. Em cada nó de uma árvore, o valor da divisão é sorteado entre o mínimo e o máximo das amostras de treino que chegam ao nó \cite{liu2008}, e uma janela além do valor extremo do treino percorre o mesmo caminho em qualquer distância além dele. Na AP6, com $d_1$, $d_2$ e $d_5$ além da faixa do treino e os demais descritores na mediana, o escore ficou abaixo do limiar nos modelos da guilhotina, da calandra, do compressor e da base pública, e com os seis descritores fora passou dele em todos os modelos (seção~\ref{sec:res-ap6}). O detector marca uma janela que se afasta do treino em vários descritores ao mesmo tempo e não marca uma mudança restrita aos descritores que a assinatura de desbalanceamento altera, por maior que ela seja. Na base pública, as janelas de 4E tinham, em mediana, os seis descritores fora da faixa, e o escore delas ficou entre \descEscoreCincoPctQuatro\ e \descEscoreMaxQuatro, logo acima do limiar de \descLimiarQuatro. Com outro treino, nenhuma delas foi marcada em \descReplicasSemMarcacaoQuatro\ das \apTresReplicas\ réplicas (seção~\ref{sec:res-ap3}).

O treino de cada máquina teve de \campoJanelasGuilhotinaUm\ a \campoJanelasCalandraUm\ janelas de um mesmo ensaio, uma a cada \campoCicloUm~s, e não cobriu a variação ao longo de semanas, com carga, temperatura e nova fixação do sensor. O falso alarme ocorreu entre rodadas da mesma máquina separadas por \campoIntervaloCalandra\ e entre duas montagens da mesma bancada. Em laboratório, \citeonline[p.~330--331]{manson2000} observaram resultado semelhante. Com a condição normal de um período curto, os dados de uma estrutura sem dano passaram do limiar quando a temperatura da sala subiu de \SI{20}{\celsius} para \SI{31}{\celsius}. Pelo axioma IVb de \citeonline[p.~1640]{worden2007}, sem uma extração de descritores feita para isso, uma medida mais sensível ao dano é também mais sensível às mudanças das condições de operação e de ambiente. Os resultados indicam que a linha de base deve incluir mais de uma condição de montagem, e os dados não indicam quantas rodadas ou quais condições bastariam.

\section{Sensibilidade diante do critério da norma}\label{sec:disc-sensibilidade}

A menor mudança detectável dependeu da variação de $d_1$ entre janelas do sinal medido, que reúne a vibração da máquina e a resposta da fixação do transmissor. Pela mesma regra, ela ficou de \apCincoRazaoMin\ a \apCincoRazaoMax\ vezes a do sensor em repouso, de \apCincoMmd~mm/s (seção~\ref{sec:res-ap5}). Se a variação do repouso e a das máquinas forem independentes, elas se somam em quadratura, e, mesmo no ensaio com a menor razão, eliminar toda a variação do repouso reduziria a menor mudança detectável em cerca de \SI{5}{\percent}. Um sensor com menos ruído teria, portanto, pouco efeito sobre a menor mudança detectável desses ensaios, e o que resta da variação é da máquina e da fixação, que os dados não separam.

A comparação com \SI{0,7}{\milli\meter\per\second} supõe que a fixação transmite a vibração da carcaça sem ganho nem atenuação, o que não foi verificado. A menor mudança detectável relativa à mediana de $d_1$, de \hDoisRelDobradeiraUm~\% na dobradeira a \hDoisRelGuilhotinaDois~\% na segunda rodada da guilhotina (Tabela~\ref{tab:h2}), não depende de um ganho constante da fixação, porque esse ganho multiplica a mediana e o quantil pelo mesmo fator. Em valor absoluto, a menor mudança detectável da segunda rodada da calandra só ficaria abaixo de \SI{0,7}{\milli\meter\per\second} na carcaça com um ganho da fixação maior que \num{4,6}, razão entre \hDoisMmdCalandraDois\ e \SI{0,7}{\milli\meter\per\second}. As decisões de H2 e H3 também se mantêm para qualquer referência de mudança até \SI{2,8}{\milli\meter\per\second}, porque a menor mudança detectável da segunda rodada da calandra passou desse valor e nenhuma máquina chegou a \SI{50}{\percent} de janelas marcadas até ele. A referência de \SI{0,7}{\milli\meter\per\second} foi derivada do critério de mudança da série ISO 10816 descrito por \citeonline{schwirzer2002}, e o texto da ISO 20816-3 não foi conferido (seção~\ref{sec:hipoteses}).

A mediana de várias janelas por medição não levou a menor mudança detectável abaixo de \SI{0,7}{\milli\meter\per\second} nos ensaios com mais variação. Com até cinco janelas, cobrindo cerca de um minuto, ela continuou acima desse valor na calandra e no compressor (seção~\ref{sec:res-ap2}). No compressor, a autocorrelação de \apDoisAutocorrCompressorUm\ entre janelas consecutivas indica uma variação lenta, que uma mediana de janelas vizinhas reduz pouco. A reamostragem da AP4 manteve esse quadro, com todas as réplicas da calandra e do compressor acima de \SI{0,7}{\milli\meter\per\second} e as da dobradeira abaixo (seção~\ref{sec:res-ap4}). Na segunda rodada da guilhotina, \apQuatroAcimaGuilhotinaDois~\% das réplicas passaram desse valor, e a posição dela em relação ao critério é incerta. Separar a variação lenta de uma falha exigiria conhecer as condições de que ela depende, como a carga ou o ciclo de trabalho, que a coleta não registrou.

Nenhum detector marcou metade das janelas com uma assinatura de \SI{0,7}{\milli\meter\per\second}. Com a fase sorteada de H3, o detector principal não chegou a \SI{50}{\percent} até \SI{2,8}{\milli\meter\per\second}, quatro vezes o critério, e com a fase fixa da AP1 também não (seção~\ref{sec:res-ap1}). Na guilhotina e no compressor, a assinatura levou $d_5$ além da faixa do treino, e o escore máximo com os descritores que ela altera ficou abaixo do limiar, de modo que o limite está no escore do detector (seção~\ref{sec:res-ap6}). Na calandra, a componente na rotação já presente, com mediana de \hTresBaseUmXCalandra~mm/s, fez a assinatura mudar pouco $d_5$. Na dobradeira, o ponto com os três descritores fora da faixa passa do limiar, mas a assinatura de \SI{2,8}{\milli\meter\per\second} levou $d_5$ além da faixa em \apSeisJanelasDcincoAcimaDobradeira\ das \hTresTesteDobradeira\ janelas de teste. A regra por mediana e MAD, que avalia cada descritor separadamente, chegou a \SI{50}{\percent} com a fase fixa na dobradeira, com \apUmMadAcinquentaZeroDobradeira~mm/s, e na guilhotina, com \apUmMadAcinquentaZeroGuilhotina~mm/s.

As medianas axiais de \SIrange{10}{184}{\hertz} da dobradeira, da calandra e do compressor ficaram acima do limite entre as zonas C e D do grupo 2 em fundação rígida, e a direção axial teve a maior velocidade nessas três máquinas (seção~\ref{sec:res-base}). O operador, porém, descreveu as três máquinas como saudáveis. Os dados não separam quatro explicações: a resposta da fixação, com fita dupla face, invólucro e espumas entre a carcaça e o sensor (seção~\ref{sec:ensaios}); o ponto de medição na carcaça do motor; a massa do transmissor; ou máquinas de fato fora da zona B. Nenhuma delas pode ser descartada sem um sensor de referência no mesmo ponto.

\section{Falhas observáveis na banda de 184~Hz}\label{sec:disc-banda}

A banda de \SIrange{10}{184}{\hertz} contém as frequências associadas ao desbalanceamento e ao desalinhamento nas máquinas da coleta. Com rotação perto de \SI{30}{\hertz} na dobradeira, na guilhotina e na calandra e perto de \SI{60}{\hertz} no compressor, a componente na rotação e a do dobro da rotação estão dentro da banda (seção~\ref{sec:res-base}). O defeito localizado de rolamento é observado por outro mecanismo. Os impactos excitam ressonâncias estruturais, e a taxa de repetição do defeito aparece como modulação dessas ressonâncias, que a análise de envelope extrai numa banda alta \cite{randall2011}. Com o filtro interno de banda nominal de \SI{184}{\hertz}, o transmissor deste trabalho corresponde ao acelerômetro de baixa taxa do nó de \citeonline{jakobsen2024}, usado só para desbalanceamento e rotação (seção~\ref{sec:fund-rolamento}).

Na guilhotina, \picosRebatidos\ linhas do espectro coincidem com ordens acima de \SI{500}{\hertz} rebatidas pela amostragem, o que é compatível com conteúdo de alta frequência que o filtro interno, digital, não removeu (seção~\ref{sec:res-base}). Conteúdo rebatido entra na curtose e no fator de crista calculados na banda, e um detector pode marcar uma janela por esse conteúdo sem relação com a falha que se quer detectar.

A fração marcada também não indica a qual mudança o detector respondeu. Na base pública, o nível 1 foi marcado em \hQuatroFracaoUm~\% das janelas com a mediana de $d_5$ próxima da de 0E e com curtose e fator de crista mais altos (seção~\ref{sec:res-h4}). Do nível 2 ao 4, a mediana de $d_5$ cresceu com o desbalanceamento, de \hQuatroDcincoMedianaDois\ a \hQuatroDcincoMedianaQuatro, e a fração marcada não seguiu essa ordem, com \hQuatroFracaoTres~\% no nível 3. Com o treino reamostrado, a fração marcada nesses níveis foi de quase nenhuma janela a todas (seção~\ref{sec:res-ap3}). Nesses dados, a componente na rotação acompanhou a severidade do desbalanceamento melhor que a decisão do detector. O trabalho detecta afastamento da operação normal e não diagnostica o tipo nem a severidade da falha.

\section{O sistema de aquisição diante dos requisitos}\label{sec:disc-sistema}

A cadeia de aquisição entregou os dados de que a avaliação dependia. Nos trechos analisados, as \campoJanelas\ janelas chegaram com os \num{26} fragmentos, a sequência das \campoLeituras\ telemetrias \campoLacunasTexto, e o resumo SHA-256 de cada janela foi conferido na exportação (seção~\ref{sec:res-aquisicao}). O número de sequência, a contagem de inicializações e a causa do reinício, que a telemetria leva para separar uma perda no ar de um reinício (seção~\ref{sec:sis-transmissor}), identificaram os reinícios por queda de tensão da bateria e as rajadas incompletas desse período, que ficaram fora dos trechos analisados.

A transmissão do sinal bruto permitiu examinar as refutações com os mesmos dados. As seis análises posteriores recalcularam a injeção da assinatura, a menor mudança detectável, o falso alarme com outras sementes e o escore do detector sobre as janelas gravadas (seção~\ref{sec:res-posteriores}). Em H4, o mesmo detector foi aplicado a uma base pública, com o sinal limitado à banda do transmissor. Num nó que executa o detector no microcontrolador e envia só a decisão ou os indicadores, uma refutação como a de H3 não poderia ser examinada com outra regra sobre as mesmas medições, porque a janela não fica guardada. Os indicadores calculados no transmissor seguem na telemetria deste trabalho, mas não entraram na avaliação.

O custo do envio do sinal bruto é o tempo de rádio. A rajada completa durou \rajadasDuracaoMediana~s, mediana na coleta (seção~\ref{sec:res-aquisicao}), contra \SI{0,411}{\second} da telemetria, e ocupou de metade a dois terços dos ciclos de \campoCicloUm\ e de \campoCicloDois~s. Com esses ciclos, o transmissor não entrou em sono profundo (seção~\ref{sec:sis-transmissor}). O consumo não foi medido, e os dados não permitem estimar a autonomia de um transmissor que envie a janela em todo ciclo. A configuração prevê janela a cada $N$ ciclos e rotina de horários, que reduzem o tempo de rádio por hora e não foram avaliadas na coleta.

O reenvio dos fragmentos ausentes completou a maior parte das rajadas fora do período de reinícios. Fora desse período, \rajadasForaCompletas\ das \rajadasFora\ rajadas chegaram completas, e das \rajadasForaIncompletas\ incompletas, \rajadasForaSemFragmento\ não tiveram nenhum fragmento recebido, o que o reenvio não corrige (seção~\ref{sec:res-aquisicao}). Das \rajadasFalhasTotal\ rajadas incompletas do dia, \rajadasFalhasReinicio\ ocorreram durante os reinícios por bateria. Com RSSI abaixo de \SI{-110}{dBm}, \rajadasAbaixoCentoDezCompletas\ das \rajadasAbaixoCentoDez\ rajadas chegou completa. Uma janela incompleta chega ao banco com a lista dos fragmentos ausentes, preenchidos com zero, e pode ser excluída da análise sem que a perda passe despercebida. Nenhuma janela dos trechos analisados ficou incompleta.

A verificação cobre melhor a lógica do protocolo e do gravador que o código ligado ao \textit{hardware}. Dos \sisTestesCasos\ casos de teste do \textit{firmware}, \sisTestesCampoCasos\ se aplicam a funções dos \textit{firmwares} da coleta, e o acesso ao sensor, ao rádio e à rede não tem teste no computador (seção~\ref{sec:sis-testes}). Três dos cinco defeitos achados nos ensaios de 2 de novembro estavam no código de rede do \textit{gateway}, que não tem teste no computador (seção~\ref{sec:sis-defeitos}). A conformidade com a norma de codificação não foi conferida por ferramenta, e a análise de modos de falha não foi refeita para o \textit{gateway}.

A versão da coleta não atende a três necessidades de um nó instalado por meses: a atualização sem cabo do transmissor, um gravador no transmissor para as janelas medidas sem o \textit{gateway} e a autenticação dos comandos pelo rádio (seção~\ref{sec:sis-limites}). O retorno automático da configuração (R5) não foi ensaiado no \textit{hardware}, e o canal cifrado (R8) só vale depois da configuração inicial. A coleta de um dia, com o operador presente, não dependeu dessas funções, que passam a ser requisitos quando o transmissor fica instalado sem acompanhamento.

\section{Relação com os trabalhos publicados}\label{sec:disc-literatura}

Os resultados reproduzem em campo o efeito de uma condição normal curta descrito em laboratório por \citeonline{manson2000} e apoiam o axioma IVb de \citeonline{worden2007} (seção~\ref{sec:disc-falso-alarme}). Eles também acrescentam aos nós de baixo custo publicados dois limites do detector que os trabalhos consultados não informam, a dependência do falso alarme em relação à semente e ao treino e o teto do escore fora da faixa do treino.

Nas duas implementações embarcadas do Isolation Forest que informam o falso alarme ou permitem calculá-lo, ele foi de \SI{22}{\percent} a \SI{38}{\percent} da operação normal, em conjuntos de dados próprios \cite{chevtchenko2023,kolok2025}. O falso alarme da calandra, de \hUmCalandraIFDireto~\% com a semente registrada, e o da base pública, de \hQuatroFracaoZero~\%, ficaram na mesma ordem de grandeza ou acima, com máquinas, conjuntos e limiares diferentes. Com \campoJanelasGuilhotinaUm\ a \campoJanelasCalandraUm\ janelas de treino, porém, o falso alarme da calandra variou de \apTresTreinoInfCalandra~\% a \apTresTreinoSupCalandra~\% com o treino reamostrado (seção~\ref{sec:res-ap3}). Um falso alarme publicado sem o tamanho do treino e sem a variação com a semente não indica se ele se repete numa nova linha de base.

\citeonline{antonini2023} treinaram e executaram o Isolation Forest no próprio ESP32 e deixaram a acurácia fora do escopo. O teto do escore e a dependência da semente medidos aqui decorrem do detector e do tamanho do treino, e não do lugar em que o detector é executado, de modo que uma implementação embarcada com treino do mesmo tamanho estaria sujeita aos mesmos limites. Com o detector no microcontrolador e só a decisão transmitida, as análises posteriores deste trabalho não seriam possíveis (seção~\ref{sec:disc-sistema}).

\citeonline{kolok2025} avaliaram o detector sem integração para velocidade e sem critério em \si{\milli\meter\per\second}. A comparação com o critério de mudança da norma, feita aqui, acrescenta uma medida que o falso alarme não fornece. Na calandra, com falso alarme abaixo dos \SI{20}{\percent} de refutação de H1, a menor mudança detectável passou de \SI{0,7}{\milli\meter\per\second} nas duas rodadas (seção~\ref{sec:res-h2}).

Os estudos de adequação de MEMS ao monitoramento de condição testaram sensores com banda de \SI{1500}{\hertz} ou mais \cite{albarbar2008,fidali2024}. Na banda de \SIrange{10}{184}{\hertz}, a menor mudança detectável do sensor de consumo em repouso foi de \apCincoRazaoMin\ a \apCincoRazaoMax\ vezes menor que a das máquinas (seção~\ref{sec:res-ap5}), e o ruído do sensor não limitou a detecção nessa banda. O defeito de rolamento, que exige banda alta, continua fora do alcance desse sensor, como no papel que \citeonline{jakobsen2024} atribui ao acelerômetro de baixa taxa (seção~\ref{sec:disc-banda}).

Em sensores da indústria pesada, \citeonline{hermansa2022} reduziram em média \SI{90,25}{\percent} a medida de falsos alarmes de um detector de agrupamento por densidade, com uma adaptação do modelo e um modelo de correção (seção~\ref{sec:fund-limiar}). Os resultados deste trabalho indicam a necessidade de um mecanismo desse tipo, ou de uma linha de base com mais de uma condição de montagem, e nenhum dos dois foi avaliado aqui.

\section{Ameaças à validade}\label{sec:limitacoes}

As ameaças à validade estão agrupadas nos quatro aspectos descritos por \citeonline[p.~153--154]{runeson2009}. A validade de construto trata de quanto as medidas representam o que a pesquisa pretende medir, e a validade interna, de fatores não controlados quando se examina uma relação de causa. A validade externa trata de até onde os resultados se generalizam, e a confiabilidade, de quanto os dados e a análise dependem de quem os produziu. Os procedimentos adotados para reduzir essas ameaças estão na seção~\ref{sec:validade-proc}.

Na validade de construto, três medidas representam só em parte o que as hipóteses tratam. A velocidade eficaz de \SIrange{10}{184}{\hertz} é um limite inferior da velocidade de banda larga da norma, e o valor em \si{\milli\meter\per\second} foi medido depois da fita dupla face, do invólucro, das espumas e das fitas do bloco eletrônico, cuja resposta não foi medida. A comparação com \SI{0,7}{\milli\meter\per\second} supõe que essa fixação transmite a vibração sem ganho nem atenuação (seção~\ref{sec:disc-sensibilidade}). A assinatura de H3 e da AP1 soma às janelas reais só a componente na rotação de um desbalanceamento, sem a resposta da máquina a uma massa real. Em H4, a bancada tem motor pequeno, eixo de \SI{12}{\milli\meter}, nenhuma carga e sensores uniaxiais em dois pontos, e o sinal foi limitado à banda do transmissor sem passar pelo sensor dele.

Na validade interna, a principal ameaça são as mudanças simultâneas entre as rodadas da mesma máquina. A fixação, a alimentação, o intervalo entre ciclos e a temperatura do sensor mudaram ao mesmo tempo, na guilhotina o sensor ficou invertido, e a orientação dos eixos não foi padronizada entre as máquinas (seção~\ref{sec:ensaios}). Essas condições não foram controladas, e os dados não atribuem o falso alarme a nenhuma delas. Parte dos dados também foi vista antes do registro. Os limites de falso alarme de H1 foram fixados depois das análises exploratórias (seção~\ref{sec:preliminares}), a AP1 e a AP2 foram desenhadas depois dos resultados de H2 e H3, a AP3 a AP6, depois de todos os resultados, e a descrição das janelas do nível 4 foi acrescentada fora do registro. Os resultados dessas análises não têm o peso de uma hipótese registrada antes de qualquer cálculo.

Na validade externa, a coleta cobriu uma instalação, um dia e no máximo duas rodadas por máquina, separadas por \campoIntervaloCalandra\ na calandra e por \campoIntervaloGuilhotina\ na guilhotina, com máquinas saudáveis segundo o operador e sem carga. Os resultados valem para o Isolation Forest e a regra por mediana e MAD com linhas de base de \campoJanelasGuilhotinaUm\ a \campoJanelasCalandraUm\ janelas de uma rodada. Eles não se estendem a outros detectores, a linhas de base mais longas nem a outras instalações. O enlace foi avaliado sem medir a distância entre o transmissor e o \textit{gateway}, e a antena do \textit{gateway}, feita para \SI{2,4}{\giga\hertz} e \SI{5}{\giga\hertz}, foi usada em \SI{433}{\mega\hertz}. Por isso, as rajadas com sinal abaixo de \SI{-110}{dBm} não podem ser atribuídas à distância nem à antena separadamente (Figura~\ref{fig:enlace}).

Na confiabilidade, a principal ameaça é o tamanho do treino. Com \campoJanelasGuilhotinaUm\ a \campoJanelasCalandraUm\ janelas, o falso alarme dependeu da semente do gerador de números aleatórios do detector a ponto de mudar a classificação de H1, e o veredito registrado vale para a semente 0 (seção~\ref{sec:res-ap3}). Os intervalos de quantis extremos, como o da menor mudança detectável, são aproximados com amostras desse tamanho (seção~\ref{sec:res-ap4}). As demais dependências do pesquisador foram reduzidas pelos procedimentos da seção~\ref{sec:validade-proc}, e os dados e os programas permitem recalcular todos os números do texto.
